\section{Experimental apparatus}
\label{sec:experiment}

This section summarizes the hardware used to produce and record the alpha-tagged recoil events
analyzed in this work. The current apparatus is housed within the AXOLOTEL-0 system, a cryogenic
and gas handling platform repurposed from the WiLiX setup \cite{arazi_first_2013, arazi_first_2015,
erdal_e_liquid_2015}.

\begin{figure}[t]
    \centering
    \includegraphics[width=0.95\textwidth]{setup_photos.png}
    \caption{Photograph of the inner detector assembly showing the grounded cathode plate,
        THGEM amplification stage, and PMT arrangement.}
    \label{fig:setup_photo}
\end{figure}

\subsection{Chamber and gas system}
While the AXOLOTEL-0 system is fully capable of liquid noble operations, for this experimental
campaign it is operated strictly with pure xenon gas at a pressure of 1 bar. Prior to data
acquisition, the gas is circulated through the purification line for 2-3 days, utilizing a hot
getter to systematically remove electronegative impurities. The electron drift time is continuously
monitored during this period until complete stabilization is achieved. For further details on the
gas line and thermodynamic control, we refer the reader to the foundational literature
\cite{arazi_first_2013, arazi_first_2015,}.

\begin{figure}[t]
    \centering
    \includegraphics[width=0.95\textwidth]{setup_schematic.png}
    \caption{Schematic of the internal detector geometry.
        Alpha decays on the silicon detector produce a prompt signal and
        eject a recoiling daughter ion into the highly constrained 4.8~mm drift region.}
    \label{fig:setup_schematic}
\end{figure}

\subsection{Detector geometry and active source}
The internal detector assembly operates as a highly constrained drift chamber
(Figure~\ref{fig:setup_photo} and Figure~\ref{fig:setup_schematic}). The main component in the
setup is the active source, which consists of a grounded flat plate housing a custom-built central
silicon PIN detector. The active surface of the silicon is coated with an extremely diluted,
ultrapure \isotope{Th}{228} layer with an activity of 50~Bq.

A homogeneous electric drift field is established across a narrow 0.48~cm gap between the cathode
plate and the amplification stage. The proportional electroluminescence (EL) region is driven by a
0.04~cm Thick Gas Electron Multiplier (THGEM). A Hamamatsu R8520-406-001 photomultiplier tube (PMT)
is positioned 20~mm from the THGEM anode surface to efficiently capture the secondary scintillation
(S2) light.

When an alpha decay occurs, the alpha particle deposits its energy in the silicon, producing a
prompt signal, while the recoiling daughter ion (e.g., \isotope{Ra}{224}) is ejected backward into
the xenon gas volume.

\subsection{Data acquisition and readout}
The fast optical readout from the PMT and the prompt energy signals from the silicon PIN detector
are recorded using a Tektronix DPO 5034 oscilloscope operating at a multi-channel sampling rate of
5~GS/s. Events are collected continuously across the full range of alpha energies. The selection of
specific isotope decays is performed entirely offline by setting strict limits on the energy range
of each alpha peak, allowing for a comprehensive, multi-isotope analysis from a single dataset.

The code employed for data acquisition and analysis is available in the open RaTag repository. The
software utilizes a flat, functional architecture driven by per-run, self-documenting YAML
configuration files to ensure complete analytical reproducibility without hidden states.